\section{The local annular remainder}
\label{sec:packet-annulus}
This section proves the annular estimate used in
Section~\ref{sec:packet-estimates}. The comparison is made before
Gaussian integration, on the actual fine configuration, coarse
configuration and predictor tags. In particular, removing a remote
cutoff does not permit us to replace the predictor by a newly chosen
one. The argument needs only the local conditions stated below.

Fix the block length $L$, the predictor patch, the positive direct
coefficients $a_f,a_c$, and the predictor penalty $p>0$. Retain the
exponent convention of Section~\ref{sec:packet-estimates}, with large
parameter $h=h(x)$. Put $t=P(h)/\sqrt h$, let $s$ be the coarse
threshold, and let $B$ and $R$ be the shell and truncation radii.
We use $s\le2t$, $R=\lfloor B/100000\rfloor$, and the eventual
geometric inequalities
\[
 100(1+B)\le c,\qquad
 10000(\rho+R_{\rm macro}+L+10)\le B,
\]
where $c$ is the coarse period and $\rho$ the predictor reach.
Both $B$ and $R$ have upper exponent zero, whereas
$e^{-\sigma R}$ has every negative upper exponent for fixed
$\sigma>0$. These properties follow from $B\asymp(\log x)^2$.
All thresholds in this section are chosen after the fixed local data;
all estimates are uniform in $c$.

The incoming linear data are the path and kernel bounds of
Section~\ref{sec:seed-guides}, the flat identities of
\eqref{eq:unhit-flat-stack}, and the additional orthogonality
\begin{equation}\label{eq:annulus-null-row}
 N_0B_cd_c=0.
\end{equation}
The last identity is the null-row part of the same flat decomposition.
Its construction, like that of the flat identities, belongs to the
flat-stack input. We will identify exactly where it is used.

\subsection{A common local chart}
Call a coarse patch clean when no coarse flag occurs there. On a
clean patch every plaquette chord is at most $s$. The axial gauge
argument in Section~\ref{sec:seed-tree-patch}, applied on a box of
radius $d+1$ with $d\le32B$, gives a gauge in which the bond chords
on the smaller box are bounded by
\[
 u_0=16\{2(32B+1)+1\}^{5}N_{\rm plaq}(2t).
\]
Thus $u_0$ has upper exponent $-1/2$. For large $x$ these bonds
are in the logarithm chart. Take their logarithms and set all other
bond logarithms to zero. This gives a full array $b$ with
$\|b\|_\infty\le C u_0$ whose exponential agrees with the gauged
coarse configuration on the required patch. No assertion about the
coarse configuration outside that patch is needed.

Write $V^\uparrow$ for the block replica of $V$. Suppose, on a
patch about $o$, the fine configuration and a colour array $u$
satisfy
\begin{equation}\label{eq:annulus-graft}
 U_e=\exp(u_e)V^\uparrow_e,\qquad \|u_e\|\le w(x),
\end{equation}
where $w\ge0$ has upper exponent $-1/2$. Under the block-constant
lift of the coarse gauge, $u_e$ is conjugated at its initial point;
its norm is unchanged. Since a replica bond is either the identity
or one coarse bond, the chord triangle inequality gives
\[
 r(U_e)\le Cw+Cu_0.
\]
Consequently both fields have small bonds in this same gauge, with
a bound $H(x)\ge w(x)$ of upper exponent $-1/2$. The coarse
logarithm extension also has norm at most $H$ after enlarging $H$.

For any prescribed $\delta>0$, all predictor reads on a sufficiently
small such patch lie in the minimizing branch and have chord at
most $\delta$. Here is the uniformity argument. At the identity
configuration the normalized minimizing frames and the best
connector are the identity. Their continuity was proved in
Section~\ref{sec:predictor-branch}. The branch inequalities are
strict there. The finitely many frame tags, connector tags and axes
therefore have a common neighbourhood on which these inequalities
hold and the predictor chord is less than $\delta$.
In the product topology on cover configurations, this neighbourhood
contains a cylinder restricting finitely many bonds. Enclose these
bonds in one fixed patch and choose one positive chord radius for
all of them. Every torus configuration whose lifted patch satisfies
these restrictions belongs to the same neighbourhood, independently
of its period. For large $x$ the patch fits within $B/2$ and its
bonds satisfy the restrictions. In particular, the normalized tag
likelihoods are exactly one there: each tag has probability the
reciprocal of the number of tags in the minimizing branch.
The four-link triangle inequality also makes every touched fine
plaquette chord at most $\delta$. These statements hold for the
original tags, including tags retained after remote masks were erased.

\subsection{Affine approximation at the actual coarse field}
Let $\mathcal C(U,V)$ be the row field consisting of the two fine
curvature rows and the predictor row:
\[
 \mathcal C(U,V)=\bigl(X_f(U),\sqrt p\,F(VQ(U)^{-1})\bigr).
\]
Here $X_f$ includes the truncated path row and the direct curvature
row, and $Q$ is the actual predictor. Each row norm is the physical
colour norm at one address; an array norm is its supremum over
addresses. No whole-volume Euclidean norm is used in these estimates. Its full linearization is
\[
 \mathcal C^0(a,b)
   =\bigl(B_fd_fa,\sqrt p\,(b-\mathcal Qa)\bigr).
\]
The chart estimates of Section~\ref{sec:seed-unhit-comparison}
give, for $\|a\|,\|b\|\le H$,
\begin{equation}\label{eq:annulus-origin-error}
 \|\mathcal C^0(a,b)\|_\infty\le C H,\qquad
 \|\mathcal C(e^a,e^b)-\mathcal C^0(a,b)\|_\infty
 \le C H\{(1+R)H+e^{-\sigma R}\}.
\end{equation}
The predictor term has a quadratic error; the path row has the
displayed additional truncation error. The right side has upper
exponent $-1$.

Define $A_\infty a=\mathcal C^0(a,0)$, and let $A_R$ be its
truncated version. Absolute exponential summability of the path
coefficients and the four-bond curl estimate give
\begin{equation}\label{eq:annulus-linear-rows}
 \|A_\infty a\|_\infty+\|A_Ra\|_\infty\le C\|a\|_\infty,
 \qquad
 \|(A_\infty-A_R)a\|_\infty
 \le C\|a\|_\infty e^{-\sigma R}.
\end{equation}
Indeed, each face curl has norm at most $4\|a\|_\infty$.
For a discarded path of length $\ell>R$, insert
$e^{-\sigma R}e^{\sigma\ell}\ge1$ into its absolute coefficient
bound and sum. The predictor part is local and is not truncated.

Let $A_R^V$ be the actual transported row: input colours are
parallel transported along the prescribed coarse routes before the
same cut linear row is applied. If the coarse bonds along these
routes have chord at most $h_0$, and the routes have length at most
$d$, then
\begin{equation}\label{eq:annulus-transport-error}
 \|(A_R^V-A_R)v\|\le
 2C D_{\rm col}\,d h_0\,\|v\|_\infty.
\end{equation}
The route product has chord at most $dh_0$ by the triangle
inequality. The bound
$\|\operatorname{Ad}(g)-1\|\le2D_{\rm col}r(g)$ then bounds
each changed input, and \eqref{eq:annulus-linear-rows} bounds the
resulting row. Only inputs within distance $R+\rho+4$ occur;
$4(R+\rho+4)\le d$ ensures the routes remain in the controlled
patch. One may set all other input differences to zero to apply
the supremum-norm bound.

We claim that, under \eqref{eq:annulus-graft} on a clean patch of
radius $B/2$, there are nonnegative functions $M,E$, of upper
exponents $-1/2$ and $-1$ respectively, such that at its central row
\begin{equation}\label{eq:annulus-affine-error}
 \|\mathcal C(U,V)\|\le M,\qquad
 \|\mathcal C(U,V)-\mathcal C(V^\uparrow,V)-A_R^Vu\|\le E.
\end{equation}
To prove this, work in the common gauge above and restrict to the
smaller radius $d=\lfloor B/4\rfloor$. Extend $u$ by zero outside
this patch, obtaining $v$ with $\|v\|_\infty\le H$. Let $a$ be
the linear replica of the coarse logarithm extension $b$; then
$\|a\|_\infty\le H$ and $e^a=(e^b)^\uparrow$.

The local product logarithm satisfies
\[
 e^{w_e}=e^{v_e}e^{a_e},\qquad
 \|w-(v+a)\|_\infty\le JH^2,\qquad
 \|w\|_\infty\le J'H.
\]
For completeness, $(v,a)\mapsto\log(e^ve^a)$ is analytic near
zero, vanishes at zero and has derivative $(v,a)\mapsto v+a$.
Integrate its second derivative on a smaller closed ball to obtain
the quadratic remainder. The first-derivative bound gives the last
inequality, bond by bond, with constants independent of the number
of bonds. Locality identifies the actual central rows with those
at $(e^w,e^b)$ and $(e^a,e^b)$.

Put $z=\mathcal C^0(w,b)$ and $z_0=\mathcal C^0(a,b)$.
Linearity gives the exact decomposition
\begin{align*}
 \mathcal C(U,V)-\mathcal C(V^\uparrow,V)-A_R^Vu
 ={}&\mathcal C(U,V)-z+A_\infty(w-v-a)\\
 &+z_0-\mathcal C(V^\uparrow,V)
 +(A_\infty-A_R)v+(A_R-A_R^V)v.
\end{align*}
The equality uses locality to replace $u$ by $v$ in the transported
row. Equations \eqref{eq:annulus-origin-error}--
\eqref{eq:annulus-transport-error} bound these five terms by
\[
 2E_0+CJH^2+CH e^{-\sigma R}+C'BH^2,
\]
where $E_0$ has upper exponent $-1$. This proves the error bound
in \eqref{eq:annulus-affine-error}; the bound on the row itself
follows from its linear bound plus $E_0$. Gauge covariance of all
three actual rows and invariance of the colour norm return the
estimate to the original gauge.

\subsection{Subtracting the selected squares}
Let $g$ be the coarse guide row and assume the already established
guide bound $\|g\|\le A t$. Define the actual and affine residuals
\[
 v=\mathcal C(U,V)-g,\qquad
 v_0=\mathcal C(V^\uparrow,V)-g+A_R^Vu.
\]
They satisfy $\|v\|\le m=M+At$ and $\|v-v_0\|\le E$.
The elementary identity for a difference of squares implies
\begin{equation}\label{eq:annulus-square-error}
 \left|\frac{\|v\|^2}{2}-\frac{\|v_0\|^2}{2}\right|
 \le Z:=\frac{E(2m+E)}2.
\end{equation}
Indeed $|\|v\|-\|v_0\||\le E$ and
$\|v_0\|\le m+E$. Hence $Z$ has upper exponent $-3/2$.

Write $\operatorname{flaw}(q)=(r(q)^2-\|F(q)\|^2)/2$.
The curvature argument in Section~\ref{sec:seed-energy} gives
\[
 0\le\operatorname{flaw}(q)\le K r(q)^4/2,
 \qquad \|F(q)\|^2\ge3r(q)^2/4
\]
for small $r(q)$. The branch and plaquette bounds above put the
required arguments in this neighbourhood. Since the appropriate
direct component of $\mathcal C$ has norm at most $M$,
\[
 r(q_f)\le2M/\sqrt{a_f},\qquad
 r(q_I)\le2M/\sqrt p.
\]
Thus their flaw contributions are bounded by
\[
 X=a_f(K/2)(2M/\sqrt{a_f})^4,
 \qquad Y=p(K/2)(2M/\sqrt p)^4,
\]
both of upper exponent $-2$.

On a kept fine row below its threshold the fine energy is exactly
\[
 a_f\operatorname{flaw}(q_f)
       +\frac12\sum_{j=0}^1\|v_{f,j}\|^2.
\]
Expand the squares to verify this: their cross terms are the
guide pairings, and their direct curvature square together with
$a_f\operatorname{flaw}$ restores $a_fr(q_f)^2/2$.
The predictor energy has, without a kept condition, the exact form
\[
 p\operatorname{flaw}(q_I)+\|v_I\|^2/2.
\]
It follows from \eqref{eq:annulus-square-error} that subtracting
the two affine fine squares costs at most $X+2Z$ in absolute
value, and subtracting the predictor square costs at most $Y+Z$.
One common nonnegative bound is $X+Y+2Z$, of upper exponent
$-3/2$. At an unselected fine or predictor address no square is
subtracted. Its corresponding displayed energy is nonnegative.
This gives the required one-sided bound there, without asserting
an absolute bound on an unselected row.

\subsection{The remaining coarse rows}
Three further terms occur in the assigned energy. First, on a
coarse plaquette with $r\le s$, the coarse counterterm is
\[
 \frac{a_c}{2}\{\|F(q)\|^2-r(q)^2\}.
\]
It is nonpositive and its absolute value is at most
$a_cKs^4/2$, of upper exponent $-2$.

Second, consider a guide kernel $K(V)$ at one central row.
On the clean patch its clipped coarse field agrees with the ordinary
truncated curvature field at every address touched by this row:
the curvature is below threshold, and the exponential kept test is
automatic outside the smaller radius where it needs checking.
Kernel locality and the local agreement $V=e^b$ then give the
exact central-row equality
\[
 [K(V)\operatorname{vision}(V)]_i=[K(e^b)X_c(e^b)]_i.
\]
This equality is taken before any supremum norm is used. Add and
subtract $K(e^b)X_c^0(b)$ to obtain
\[
 \bigl\|[K(V)\operatorname{vision}(V)-K_0X_c^0(b)]_i\bigr\|
 \le C\|X_c(e^b)-X_c^0(b)\|_\infty
       +C\omega(R,H)\|X_c^0(b)\|_\infty,
\]
where $\omega(R,H)=(1+R)H+e^{-\tau R}$. The row bounds in
\eqref{eq:unhit-kernel-near} justify this inequality. The globally
small extension $b$ permits both supremum estimates on the right;
no global bound on the original $V$ is required. The right side has
upper exponent $-1$. For the null row $K=N$, its full linear
comparator vanishes by \eqref{eq:annulus-null-row}. The null
energy, half the square of this row, consequently has upper
exponent $-2$. This is the only extra flat identity used here.

Finally, the linear and quadratic mismatch rows are precisely the
gaps estimated in \eqref{eq:unhit-global}. The common local gauge
and locality argument used there apply to the present grafted
configuration, with smallness bound $H$ of upper exponent $-1/2$.
Pairing the linear gap with a guide of size $At$ gives at most
$At\,CH\omega(R,H)$ for each of the two fine-row channels.
Pairing the quadratic gap with that guide gives at most
$At\,CAt\,\omega(R,H)/2$ per channel. Hence their combined
absolute cost is at most
\[
 2At\,CH\omega(R,H)+(At)^2C\omega(R,H),
\]
of upper exponent $-3/2$. Adding all row costs gives one
nonnegative function $r(x)$ of upper exponent $-3/2$.

\subsection{The protected annulus of an actual packet}
We verify the local hypotheses on the packet itself. Its fine
configuration is formed in a fixed tree slice. On active free
bonds it is $\exp(u_e)V^\uparrow_e$; on frozen free bonds it is
the fixed slice value $Y_e$; on tree bonds it is the identity.
Let $K>0$ be a fixed shell constant and define
\[
 \ell_K(q)=\chi\bigl((r(q)/(Kt))^2\bigr)\log q,
\]
where the smooth cutoff equals one for $r(q)\le Kt$ and vanishes
for $r(q)\ge2Kt$. The local logarithm estimate gives
$\|\ell_K(q)\|\le2\gamma Kt$ everywhere. This assertion
also holds outside the chart because the cutoff is zero there;
choose $x$ so that the closed ball $r(q)\le2Kt$ lies in the
logarithm chart. The resulting map is smooth wherever its group
argument varies smoothly, by the same two local cases.

Define the actual colour signal $\widetilde u$ to be $u$ on
active bonds, $\ell_K(Y_e(V^\uparrow_e)^{-1})$ on frozen free
bonds, and zero on tree bonds. On a weak interior patch the
fine plaquettes have chord at most $t/2$, the predictor mismatch
at most $2Wt$, and the selected frames have positive likelihood.
The tree comparison \eqref{eq:tree-shell} then gives
\[
 r(U_e(V^\uparrow_e)^{-1})
 \le Kt,\qquad K=C\,N_{\rm macro}/2+2W,
\]
where $N_{\rm macro}$ counts the fine plaquettes in the fixed
macro patch. The frozen cutoff is therefore one and its logarithm
exponentiates to the original mismatch. On tree bonds both
$U_e$ and $V^\uparrow_e$ are the identity. Thus on this patch
\[
 U_e=\exp(\widetilde u_e)V^\uparrow_e.
\]
If the active variables there obey $\|u_e\|\le w(x)$ and
$2\gamma Kt\le w(x)$, the signal obeys the same bound on all
bonds. This proves \eqref{eq:annulus-graft} for the actual signal.

Let $A$ be the packet's seed set, and consider an energy address
whose coarse anchor belongs to its $200B$ neighbourhood but not
its $2B$ neighbourhood. Retain the weak fine, predictor and coarse
flag conditions through $300B$, and the bounds on active variables
through $304B$. For a bond within $B$ of this anchor, its
$B/4$ interior patch lies within the $300B$ neighbourhood and
outside the seeds. This follows from the triangle inequality:
the outer distance is at most $201B+B/4<300B$, whereas entering
the seeds would put the anchor within $B+B/4<2B$ of them.
The retained conditions therefore give the weak interior just used.
The active-variable condition is available since $201B<304B$.
The same argument supplies a clean coarse patch at the anchor.

Signed changes of orientation preserve these conclusions, allowing
one additional bond of radius for reversed bonds. Norms are
unchanged by the associated conjugations. Finally, the fine kept
test holds on this annulus: a plaquette whose test is nontrivial
has distance at most $B/25$ by the exponential threshold bound;
the weak interior makes its chord at most $t/2\le
t e^{\operatorname{dist}/r_0}$. At all other plaquettes the
right side is at least the universal chord bound $2$.
Thus the actual fine energy equals the square expression used above,
including at threshold equality.

At every selected address the extracted affine square is the square
of the actual signal row with the same tags, orientations and coarse
field. At every unselected address it is zero. We have therefore
proved, on the protected annulus,
\begin{equation}\label{eq:annulus-final}
 \text{actual assigned energy at the address}\ \ge-r(x).
\end{equation}
Here $r\ge0$ has upper exponent $-3/2$.
Summing over its at most $C B^4$ addresses per seed and averaging
over the signed channels gives exactly the annular input of
\eqref{eq:packet-assigned-energy}. This proves the value estimate;
derivative bounds for the density, covariance and chart factors
remain separate inputs to the packet jet estimate.
